\documentclass[manuscript,review,nonacm]{acmart}
%
% COMP5111 Assignment 2 -- Task 2 LLM-use report
%

\usepackage{listings}
\usepackage{xcolor}

\lstset{
  basicstyle=\ttfamily\small,
  breaklines=true,
  columns=fullflexible,
  frame=single,
  framesep=4pt,
  backgroundcolor=\color{gray!8},
  keywordstyle=\bfseries\color{blue!60!black},
  commentstyle=\itshape\color{green!50!black},
  showstringspaces=false,
}

\setcopyright{none}
\settopmatter{printacmref=false, printccs=false, printfolios=true}
\renewcommand{\footnotetextcopyrightpermission}[1]{}
\pagestyle{plain}

\title{Task 2 \textendash{} LLM-Use Report\\
       \large COMP5111 Assignment 2}
\author{Zimo Ji}
\affiliation{\institution{HKUST}\city{Hong Kong}\country{Hong Kong SAR}}

\begin{document}
\maketitle
\thispagestyle{plain}

\section{The LLM I used}
I used \textbf{ChatGPT (GPT-5.4)} through an OpenAI-compatible chat
endpoint. I did not let it run my pipeline; I pasted relevant rows of
the \texttt{*.tsv} spectra, the failing test bodies, and chunks of
\texttt{Subject.java} into the chat and asked specific questions about
suspicious Jimple statements.

\section{Where the LLM helped}

\subsection{Picking the right row when one Jimple stmt maps to two
source lines}
In \texttt{parseToken} the buggy \texttt{pos++;} (line 205) and the
harmless \texttt{pos++;} (line 208) both render as
\texttt{i2 = i2 + 1}. My TSV has two rows with that text, scored
0.617 and 0.552 (in \texttt{randoop0}), but no source line. I asked
the LLM which one was likely the buggy one. It pointed out that line
205 only fires when a terminator is hit, which is exactly the path
that fails, so it should score higher. That matched my data and let
me cite \texttt{rank=7, score=0.617} in \texttt{fault\_205.txt}
without guessing.

\subsection{One-line specs for the suspicious methods}
Reading the Javadocs of \texttt{parseToken}, \texttt{extractIntInStr},
\texttt{daysBetweenDateStrings}, and \texttt{monAbbr2month} is slow.
I asked the LLM to summarise each contract in one sentence. For
\texttt{extractIntInStr} it gave me ``return the LAST contiguous digit
run, e.g. \texttt{1234a123} $\to$ 123'', and once I had that, the
failing test
\texttt{extractIntInStr("0366-01-31") expected 31} made the bug
obvious: the code returns at the first non-digit, but the spec wants
the last digit run. I had read the function a few times and convinced
myself it ``returns the leading int'' before the LLM made me look at
the Javadoc again.

\subsection{Computing the 12 perfect-hash constants for
\texttt{monAbbr2month}}
The method dispatches on \texttt{(ch0 << 16) | (ch1 << 8) | ch2}
against 12 hard-coded integers. I asked the LLM to compute the hashes
for ``Jan'' through ``Dec'' and check them against the source. It
returned a table where every row matched except \texttt{Sep}: source
\texttt{5465466}, computed \texttt{5465456}. I confirmed by hand:
\texttt{('S' << 16) | ('e' << 8) | 'p' = 5439488 + 25856 + 112 =
5465456}. That is exactly the off-by-10 fault on line 615. Doing this
without the LLM is 12 multiplications and additions, all correct,
which I would not have trusted myself to do without a typo.

\subsection{Splitting ``Calendar bug'' from ``validity bug'' in
\texttt{daysBetweenDates}}
The failing tests
\texttt{daysBetweenDateStrings("0100-10-31","0366-01-31")} etc.\ light
up both \texttt{daysBetweenDates} (lines 442--454, the Calendar code)
and \texttt{checkValidDate} (line 488). I asked the LLM where in
\texttt{daysBetweenDateStrings} the value \texttt{Integer.MIN\_VALUE}
could come from. It pointed out that \texttt{MIN\_VALUE} only comes
from the catch block, so an exception must be thrown before the
Calendar arithmetic runs. With day=31 in the failing inputs, the only
exception source is \texttt{checkValidDate} rejecting day=31 of a
31-day month. That moved my attention straight to line 488 and I
stopped looking at the Calendar code.

\section{What the LLM was not useful for}
\begin{itemize}
\item It could not improve the ranking. Ochiai already gives the
      right ordering up to ties; the LLM does not change the score
      formula.
\item It could not, on its own, separate two source lines that share
      the same Jimple text. I had to feed it the metadata file mapping
      stmts to lines.
\item It kept proposing rewrites for \texttt{extractIntInStr} (e.g.\
      ``use a regex''), which violates the single-line patch rule. I
      had to remind it of the constraint.
\end{itemize}

\section{Net}
I used the LLM as a Javadoc reader and an arithmetic checker on top
of my Ochiai shortlist. That cut the manual inspection from ``every
row in top-30'' to ``the four lines I am submitting as faults''.
After applying my four single-line patches, all 1371 tests in the
four provided suites pass.

\noindent\textit{(Word count, excluding code: about 470.)}

\end{document}
